\section*{Chapter \ref{ch:g11:scal} --- The Scalar Product in the Plane}

\begin{solution}{exo:g11:scal:1}
$(2,5)\cdot(3,-1) = 6 - 5 = 1$.

$(1,-4)\cdot(8,2) = 8 - 8 = 0$ (orthogonal vectors).

$\vec u\cdot\vec v = 3 \times 4 \times \cos\frac\pi3 = 12 \times \frac12
= 6$.
\end{solution}

\begin{solution}{exo:g11:scal:2}
$\vect{AB}\cdot\vect{AC} = 6 \times 6 \times \cos 60^\circ = 18$.

The vectors $\vect{AB}$ and $\vect{BC}$ make an angle of
$180^\circ - 60^\circ = 120^\circ$ (place them tail to tail at $B$), so
$\vect{AB}\cdot\vect{BC} = 36 \cos 120^\circ = -18$.
\end{solution}

\begin{solution}{exo:g11:scal:3}
$\vec u \cdot \vec v = 3(t - 8) + t = 4t - 24 = 0$ gives $t = 6$.

$\vec u - \vec v = (t - 4,\ 0)$, so
$\vec u \cdot (\vec u - \vec v) = t(t-4) + 0 = t^2 - 4t$, which vanishes
for $t = 0$ or $t = 4$.
\end{solution}

\begin{solution}{exo:g11:scal:4}
$\vec u \cdot \vec v = 3 + 2 = 5$, $\norm{\vec u} = \sqrt5$,
$\norm{\vec v} = \sqrt{10}$, so
\[
\cos\theta = \frac{5}{\sqrt5\,\sqrt{10}} = \frac{5}{5\sqrt2}
= \frac{\sqrt2}{2},
\]
hence $\theta = 45^\circ$ exactly.
\end{solution}

\begin{solution}{exo:g11:scal:5}
$a^2 = 16 + 36 - 2 \times 4 \times 6 \times \cos\frac{2\pi}{3}
= 52 - 48 \times \left(-\frac12\right) = 76$, so $a = \sqrt{76} =
2\sqrt{19} \approx 8.7$.

For the second triangle, the law of cosines solved for the angle gives
\[
\cos\widehat A = \frac{b^2 + c^2 - a^2}{2bc}
= \frac{25 + 16 - 49}{2 \times 5 \times 4} = -\frac{8}{40} = -\frac15,
\]
negative: $\widehat A$ is obtuse.
\end{solution}

\begin{solution}{exo:g11:scal:6}
$\vect{AB}\,(4, -2)$ and $\vect{AC}\,(2, 4)$:
$\vect{AB}\cdot\vect{AC} = 8 - 8 = 0$, so the angle at $A$ is right.
Sides: $AB = \sqrt{16 + 4} = \sqrt{20}$, $AC = \sqrt{4 + 16} =
\sqrt{20}$, and $\vect{BC}\,(-2, 6)$ gives $BC = \sqrt{40}$. Consistency:
$BC^2 = 40 = 20 + 20 = AB^2 + AC^2$, the law of cosines with
$\cos\widehat A = 0$ (Pythagoras).
\end{solution}

\begin{solution}{exo:g11:scal:7}
Center--radius form: $(x + 2)^2 + (y - 3)^2 = 25$.

For the second circle, complete the squares:
\[
x^2 - 6x + y^2 + 2y = 6
\iff (x - 3)^2 - 9 + (y + 1)^2 - 1 = 6
\iff (x-3)^2 + (y+1)^2 = 16 :
\]
center $(3, -1)$, radius $4$.
\end{solution}

\begin{solution}{exo:g11:scal:8}
$M(x,y)$ is on the circle when $\vect{MA}\cdot\vect{MB} = 0$ with
$\vect{MA}\,(-3 - x,\ -y)$ and $\vect{MB}\,(-x,\ 2 - y)$:
\[
(-3 - x)(-x) + (-y)(2 - y) = x^2 + 3x + y^2 - 2y = 0 .
\]
For the origin: $0 + 0 + 0 - 0 = 0$, so $O$ lies on the circle ---
geometrically, $\vect{OA}\,(-3,0)$ and $\vect{OB}\,(0,2)$ are orthogonal,
so $O$ sees $[AB]$ under a right angle. (Completing the squares:
center $\left(-\frac32, 1\right)$, radius $\frac{\sqrt{13}}{2}$.)
\end{solution}

\begin{solution}{exo:g11:scal:9}
The perpendicular through $P$ has the direction vector of $d$, namely
$(-1, 2)$, as normal vector: equation $-x + 2y + c = 0$; through
$P(1,3)$: $-1 + 6 + c = 0$, $c = -5$, so $x - 2y + 5 = 0$.

Intersection with $d$: from $2x + y = 7$, $y = 7 - 2x$; substituting,
$x - 2(7 - 2x) + 5 = 5x - 9 = 0$, so $x = \frac95$, $y = \frac{17}5$:
$H\left(\frac95, \frac{17}5\right)$. Then
$\vect{PH}\left(\frac45, \frac25\right)$ and
\[
PH = \sqrt{\frac{16}{25} + \frac{4}{25}} = \frac{\sqrt{20}}{5}
= \frac{2\sqrt5}{5} \approx 0.89 .
\]
\end{solution}

\begin{solution}{exo:g11:scal:10}
Adding
\[
\norm{\vec u + \vec v}^2 = \norm{\vec u}^2 + 2\vec u\cdot\vec v +
\norm{\vec v}^2
\quad\text{and}\quad
\norm{\vec u - \vec v}^2 = \norm{\vec u}^2 - 2\vec u\cdot\vec v +
\norm{\vec v}^2,
\]
the cross terms cancel and
$\norm{\vec u + \vec v}^2 + \norm{\vec u - \vec v}^2 = 2\norm{\vec u}^2 +
2\norm{\vec v}^2$. In a parallelogram with sides $\vec u$ and $\vec v$,
the diagonals are $\vec u + \vec v$ and $\vec u - \vec v$: the sum of the
squares of the two diagonals equals the sum of the squares of the four
sides.
\end{solution}

\begin{solution}{exo:g11:scal:11}
\emph{1.} Since $I$ is the midpoint, $\vect{IB} = -\vect{IA}$ and
$IA = 2$. Then
\[
\vect{MA}\cdot\vect{MB}
= (\vect{MI} + \vect{IA})\cdot(\vect{MI} - \vect{IA})
= \norm{\vect{MI}}^2 - \norm{\vect{IA}}^2
= MI^2 - 4 .
\]

\emph{2.} The condition $\vect{MA}\cdot\vect{MB} = 12$ becomes
$MI^2 = 16$, i.e.\ $MI = 4$: the set is the circle of center $I$ and
radius $4$.
\end{solution}
